\documentclass[12pt,letterpaper]{article}
\usepackage{graphicx,textcomp}
\usepackage{natbib}
\usepackage{setspace}
\usepackage{fullpage}
\usepackage{color}
\usepackage[reqno]{amsmath}
\usepackage{amsthm}
\usepackage{fancyvrb}
\usepackage{amssymb,enumerate}
\usepackage[all]{xy}
\usepackage{endnotes}
\usepackage{lscape}
\newtheorem{com}{Comment}
\usepackage{float}
\usepackage{hyperref}
\newtheorem{lem} {Lemma}
\newtheorem{prop}{Proposition}
\newtheorem{thm}{Theorem}
\newtheorem{defn}{Definition}
\newtheorem{cor}{Corollary}
\newtheorem{obs}{Observation}
\usepackage[compact]{titlesec}
\usepackage{dcolumn}
\usepackage{tikz}
\usetikzlibrary{arrows}
\usepackage{multirow}
\usepackage{xcolor}
\newcolumntype{.}{D{.}{.}{-1}}
\newcolumntype{d}[1]{D{.}{.}{#1}}
\definecolor{light-gray}{gray}{0.65}
\usepackage{url}
\usepackage{listings}
\usepackage{color}

\definecolor{codegreen}{rgb}{0,0.6,0}
\definecolor{codegray}{rgb}{0.5,0.5,0.5}
\definecolor{codepurple}{rgb}{0.58,0,0.82}
\definecolor{backcolour}{rgb}{0.95,0.95,0.92}

\lstdefinestyle{mystyle}{
	backgroundcolor=\color{backcolour},   
	commentstyle=\color{codegreen},
	keywordstyle=\color{magenta},
	numberstyle=\tiny\color{codegray},
	stringstyle=\color{codepurple},
	basicstyle=\footnotesize,
	breakatwhitespace=false,         
	breaklines=true,                 
	captionpos=b,                    
	keepspaces=true,                 
	numbers=left,                    
	numbersep=5pt,                  
	showspaces=false,                
	showstringspaces=false,
	showtabs=false,                  
	tabsize=2
}
\lstset{style=mystyle}
\newcommand{\Sref}[1]{Section~\ref{#1}}
\newtheorem{hyp}{Hypothesis}

\title{Problem Set 2}
\date{Due: February 18, 2026}
\author{Applied Stats II}


\begin{document}
	\maketitle
	\section*{Instructions}
	\begin{itemize}
		\item Please show your work! You may lose points by simply writing in the answer. If the problem requires you to execute commands in \texttt{R}, please include the code you used to get your answers. Please also include the \texttt{.R} file that contains your code. If you are not sure if work needs to be shown for a particular problem, please ask.
		\item Your homework should be submitted electronically on GitHub in \texttt{.pdf} form.
		\item This problem set is due before 23:59 on Wednesday February 18, 2026. No late assignments will be accepted.

	\end{itemize}

	\vspace{.25cm}

\noindent We're interested in what types of international environmental agreements or policies people support (\href{https://www.pnas.org/content/110/34/13763}{Bechtel and Scheve 2013)}. So, we asked 8,500 individuals whether they support a given policy, and for each participant, we vary the (1) number of countries that participate in the international agreement and (2) sanctions for not following the agreement. \\

\noindent Load in the data labeled \texttt{climateSupport.RData} on GitHub, which contains an observational study of 8,500 observations.

\begin{itemize}
	\item
	Response variable: 
	\begin{itemize}
		\item \texttt{choice}: 1 if the individual agreed with the policy; 0 if the individual did not support the policy
	\end{itemize}
	\item
	Explanatory variables: 
	\begin{itemize}
		\item
		\texttt{countries}: Number of participating countries [20 of 192; 80 of 192; 160 of 192]
		\item
		\texttt{sanctions}: Sanctions for missing emission reduction targets [None, 5\%, 15\%, and 20\% of the monthly household costs given 2\% GDP growth]
		
	\end{itemize}
	
\end{itemize}

\newpage
\noindent Please answer the following questions:

\vspace{.25cm}

\section*{Question 1}

\begin{enumerate}
	\item
	Remember, we are interested in predicting the likelihood of an individual supporting a policy based on the number of countries participating and the possible sanctions for non-compliance.
	\begin{enumerate}
		\item [] Fit an additive model. Provide the summary output, the global null hypothesis, and $p$-value. Please describe the results and provide a conclusion.
	\end{enumerate}

\vspace{.5cm}

\noindent To fit the additive model, firstly the data was loaded. Then the contrasts were assigned to contr.treatment, which implements dummy coding and takes one category as a reference category. This ensures the model will have interpretable outputs. Finally, the glm function was used with choice as the outcome variable and countries and sanctions as explanatory variables (with no interaction term), using the logit link function. The summary of the additive model was printed.

\lstinputlisting[language=R, firstline=39, lastline=51]{PS02_KC.R}

\begin{verbatim}
	Call:
	glm(formula = choice ~ countries + sanctions, family = binomial(link = "logit"), 
	data = climateSupport)
	
	Coefficients:
	Estimate Std. Error z value Pr(>|z|)    
	(Intercept)         -0.27266    0.05360  -5.087 3.64e-07 ***
	countries80 of 192   0.33636    0.05380   6.252 4.05e-10 ***
	countries160 of 192  0.64835    0.05388  12.033  < 2e-16 ***
	sanctions5%          0.19186    0.06216   3.086  0.00203 ** 
	sanctions15%        -0.13325    0.06208  -2.146  0.03183 *  
	sanctions20%        -0.30356    0.06209  -4.889 1.01e-06 ***
	---
	Signif. codes:  0 ‘***’ 0.001 ‘**’ 0.01 ‘*’ 0.05 ‘.’ 0.1 ‘ ’ 1
	
	(Dispersion parameter for binomial family taken to be 1)
	
	Null deviance: 11783  on 8499  degrees of freedom
	Residual deviance: 11568  on 8494  degrees of freedom
	AIC: 11580
	
	Number of Fisher Scoring iterations: 4
\end{verbatim}

\noindent The global null hypothesis is that none of the input variables (countries and sanctions) in the model are significant in predicting the output variable (choice). This can be tested by running an ANOVA comparing our model to the null model.

\lstinputlisting[language=R, firstline=56, lastline=57]{PS02_KC.R}

\begin{verbatim}
	Analysis of Deviance Table
	
	Model 1: choice ~ countries + sanctions
	Model 2: choice ~ 1
	Resid. Df Resid. Dev Df Deviance  Pr(>Chi)    
	1      8494      11568                          
	2      8499      11783 -5  -215.15 < 2.2e-16 ***
	---
	Signif. codes:  0 ‘***’ 0.001 ‘**’ 0.01 ‘*’ 0.05 ‘.’ 0.1 ‘ ’ 1
\end{verbatim}

\noindent The p-value found from the ANOVA test is 2.2e-16. That is below our critical value of 0.05. So, we can reject the global null hypothesis that no variables in the model are statistically significant predictors of choice. This means that at least one variable is statistically significant.

\section*{Question 2}

	\item
	If any of the explanatory variables are statistically significant in this model, then:
	\begin{enumerate}
		
		\section*{2a}
		
		\item
		For the policy in which nearly all countries participate [160 of 192], how does increasing sanctions from 5\% to 15\% change the odds that an individual will support the policy? (Interpretation of a coefficient)
		
		\vspace{.5cm}
		
		\noindent To find the difference between the 5 percent coefficient and 15 percent coefficient, the coefficients were separately extracted from the summary of the additive model and assigned to different variables. Then the coefficients were stored in a vector because the diff function expects a vector. The diff function was used on the vector and the result was assigned to the variable coefficient\textunderscore difference, which was printed.
		
		\lstinputlisting[language=R, firstline=65, lastline=69]{PS02_KC.R}
		
		\begin{verbatim}
		[1] -0.3251028
		\end{verbatim}
		
		\noindent For the policy in which nearly all countries participate [160 of 192], increasing sanctions from 5 percent to 15 percent decreases the log-odds that an individual will support the policy by approximately 0.3251028 on average, holding all other variables constant.
		
		\lstinputlisting[language=R, firstline=75, lastline=75]{PS02_KC.R}
		
		\begin{verbatim}
		[1] 0.7224531
		\end{verbatim}
		
		\noindent For the policy in which nearly all countries participate [160 of 192], increasing sanctions from 5 percent to 15 percent decreases the odds of support for the policy by a factor of exp(-0.3251028) (approximately 0.7224531) on average.
		As 1 - 0.7224 = 0.2776, this indicates that increasing sanctions from 5 percent to 15 percent decreases the odds of support for the policy by approximately 27.76 percent.
		
		\section*{2b}
		
		\item
		For the policy in which very few countries participate [20 of 192], how does increasing sanctions from 5\% to 15\% change the odds that an individual will support the policy? (Interpretation of a coefficient)
		
		\vspace{.5cm}
		
		\noindent The steps above were repeated. The same coefficients were used and the output was identical.
		
		\lstinputlisting[language=R, firstline=84, lastline=88]{PS02_KC.R}
		
		\begin{verbatim}
		[1] -0.3251028
		\end{verbatim}
		
		\noindent For the policy in which very few countries participate [20 of 192], increasing sanctions from 5 percent to 15 percent decreases the log-odds that an individual will support the policy by approximately 0.3251028 on average, holding all other variables constant.
		
		\lstinputlisting[language=R, firstline=94, lastline=94]{PS02_KC.R}
		
		\begin{verbatim}
		[1] 0.7224531
		\end{verbatim}
		
		\noindent For the policy in which very few countries participate [20 of 192], increasing sanctions from 5 percent to 15 percent decreases the odds of support for the policy by a factor of exp(-0.3251028) (approximately 0.7224531) on average.
		As 1 - 0.7224 = 0.2776, this indicates that increasing sanctions from 5 percent to 15 percent decreases the odds of support for the policy by approximately 27.76 percent.
	    
	    \vspace{.25cm}
	    
		\noindent The interpretation of the coefficient for the policy in which very few countries participate and the interpretation of the coefficient for the policy in which nearly all countries participate are the same because this is an additive model that assumes that there is no interaction between the percent of sanctions and the number of countries participating. Changing the number of countries does not impact our interpretation of the increase in sanctions because our model assumes that the relationship between the percent of sanctions and support for the policy is the same at all levels of country participation.
		
		\vspace{.25cm}
		
		\section*{2c}
		
		\item
		What is the estimated probability that an individual will support a policy if there are 80 of 192 countries participating with no sanctions? 
	
		\vspace{.5cm}
	
		\noindent To find the estimated probability that an individual will support a policy if there are 80 of 192 countries participating with no sanctions the probability formula was used in which the \begin{math} probability = \frac{exp(\beta0 + \beta1X_{i}}{1 + exp(\beta0 + \beta1X_{i})}) \end{math}
	
		\noindent The coefficients from the additive model summary were substituted into the equation. However, as no sanctions is the reference category, there is no coefficient for it and its effect is absorbed into the intercept. Therefore, the equation is \begin{math} probability = \frac{exp(-0.27266 + 0.33636 * 1}{1 + exp(-0.27266 + 0.33636 * 1)}) \end{math}
	
		\lstinputlisting[language=R, firstline=109, lastline=110]{PS02_KC.R}
	
	\begin{verbatim}
	[1] 0.5159196
	\end{verbatim}
	
	\noindent The probability that an individual will support a policy if there are 80 of 192 countries participating with no sanctions is approximately 51.59 percent.
	
	\end{enumerate}
	
	\section*{Question 3}
	
	\item
	Would the answers to 2a and 2b potentially change if we included an interaction term in this model? Why?
	
		\begin{itemize}
		\item Perform a test to see if including an interaction is appropriate.
	\end{itemize}
	\end{enumerate}
	
	\noindent The answers to 2a and 2b would potentially change because adding an interaction term would account for the fact that the impact of an explanatory variable on the outcome variable could be influenced by another explanatory variable in the model. In this case, adding an interaction term allows for the relationship between sanctions and support for the policy to differ at different levels of participating countries and vice versa. This means that the answers for 2a and 2b would no longer always be the same.
	
	\lstinputlisting[language=R, firstline=122, lastline=126]{PS02_KC.R}
	
	 \begin{verbatim}
	 Analysis of Deviance Table
	 
	 Model 1: choice ~ countries + sanctions
	 Model 2: choice ~ countries * sanctions
	 Resid. Df Resid. Dev Df Deviance Pr(>Chi)
	 1      8494      11568                     
	 2      8488      11562  6   6.2928   0.3912
	 \end{verbatim}

\noindent The ANOVA has a p-value of 0.3912 which is larger than the critical value of 0.05, this suggest that the difference between models is not significant. For the sake of parsimony it may be better not to include an interaction effect, unless there is a theoretical backing.

\end{document}
